Tangut translation is an extreme low-resource problem: in the current repository snapshot, only 491 real Tangut-Chinese pairs are available, with 400/41/50 train-dev-test examples. The task is also structurally unusual because many targets are short, title-like Chinese strings rather than ordinary modern sentences. We evaluate three families of methods in this setting: inference-only prompting, synthetic-data supervised fine-tuning (SFT), and direct preference optimization (DPO). Structured multitask SFT is the strongest pure supervised baseline, but the strongest current learned system comes from stricter gap-filtered DPO on top of that base. With only 498 preference pairs whose reward gap is at least 0.4, the sigmoid variant reaches chrF++ 30.01, keeps 5/50 exact matches, preserves 22/27 title suffixes, and ties for the best overall score (2.22/5) under a new reference-aware judge. We further show that several apparently reasonable reference-free metrics are misaligned with the task: the gold human reference itself receives only moderate dictionary coverage and extremely poor perplexity under a modern-Chinese language model, despite achieving perfect chrF++ and perfect reference-aware judge scores. At the same time, preference optimization is not uniformly helpful: legacy DPO on a weaker base still degrades reference-aligned performance, and looser gap filtering yields adequacy gains only with substantial length or contamination costs. These findings suggest a two-stage practical lesson for ultra-low-resource Tangut short-text translation: explicit structural supervision is the right starting point, preference optimization can help only when pair quality is aggressively controlled, and evaluation must be anchored to reference-aware metrics rather than generic fluency proxies.
